% ID: MAT-ALG-INT-002
% Titolo: Numeri interi e proprietà delle potenze 2
% Disciplina: Matematica
% Area: algebra
% Argomento: numeri_interi
% Classe: Prima secondaria di secondo grado
% Difficolta: 3
% Tipo: problema_numerico
% Tag: matematica, algebra, numeri_interi
% Fonte: verifica_interi_anonimizzata
% Provenienza: materiale fornito dal docente; testo anonimizzato; soluzione redatta per Exergo

\begin{esercizio}
Disponi in ordine crescente i risultati delle seguenti espressioni:
    \[
    -\lvert -3\rvert^3;\qquad
    -(-4)^2;\qquad
    -(-2)^2\cdot(-2)^3;\qquad
    (-90)^2:(15)^2;\qquad
    \left[(-16)^4-(-2)^6(-32)^2\right]^5.
    \]
\end{esercizio}

\begin{soluzione}
I valori nell'ordine del testo sono $-27,-16,32,36,0$.
Quindi l'ordine crescente è $-27<-16<0<32<36$.
\end{soluzione}
